% ---------------------------------------------------------------------------
\section{NS-CSD: Neuro-Symbolic Continuous Solution Discovery}
\label{sec:method}
% ---------------------------------------------------------------------------

NS-CSD keeps Sentinel's problem formulation, i.e.\ the POMDP over 12 aggregate
telemetry features, the action $a_t=\langle\tau_t,f_t,d_t\rangle$ and the symbolic
shield $\mathcal{S}$, and replaces its training machinery. It drops the co-evolving
PPO attacker, the Hall-of-Fame archive, the learned critic and post-hoc
checkpoint selection. In their place it runs one continuous loop with five
interacting parts (Fig.~\ref{fig:arch}, Algorithm~\ref{alg:nscsd}):
(i) a critic-free GRPO defender trained on forked, common-random-number rollouts;
(ii) a quality-diversity \emph{Solution Archive} whose elites are heterogeneous
(neural snapshots, crystallised decision-tree programs and seed symbolic programs);
(iii) a regret-driven \emph{Scenario Curriculum} that generates attack scenarios
instead of learning an attacker; (iv) archive-guided policy improvement through
elite injection, a KL anchor and champion roll-back; and (v) certificate-gated
\emph{discovery of the symbolic shield itself}.

\subsection{Scenario space}
\label{sec:scenarios}
An attack scenario is a parameter vector
$\xi=(w,\;p_{\text{stay}},\;[i_\ell,i_h],\;[m_\ell,m_h],\;P,\;\delta,\;p_{\text{flash}})\in\Xi\subset\mathbb{R}^{14}$:
attack-family weights $w$ over \{none, SYN, UDP, HTTP, ICMP, mixed\}, the Markov
persistence of the active family, intensity and mutation ranges that are
resampled on every family switch, an optional on/off pulse (period $P$, duty
$\delta$) and a flash-crowd probability. A scenario compiles into an
open-loop attack schedule plus an \emph{exogenous tape} holding every random
quantity the simulator draws: Poisson arrivals, flash-crowd multipliers and
observation noise. Training scenarios never put weight on ICMP, so ICMP
stays held out. A coarse descriptor
$\kappa(\xi)=(\arg\max_k w_k,\;\text{intensity bin})$ splits $\Xi$ into 18 niches.

\subsection{Critic-free GRPO with common random numbers}
\label{sec:grpo}
Each iteration draws $B$ contexts $\xi_b$ from the curriculum and warms each one
up for $W\sim\mathcal{U}\{0,\dots,W_{\max}\}$ steps under the current policy,
reaching a realistic mid-episode state $s^{(b)}_0$. The context is then
\emph{forked} into a group of $G$ members that replay the same tape. That is,
they share common random numbers (CRN)~\cite{ng2000pegasus}. Member $g$ acts
for $H$ steps through the shield and collects the verifiable return
\begin{equation}
R_{b,g}=\sum_{h=0}^{H-1} r(s_h,\mathcal{S}(s_h,a_h)),\quad
r = \mathrm{SQ}-\mathrm{leak}-w_{\text{sev}}\mathbb{1}[\mathrm{SQ}<0.5]-0.25\,\mathbb{1}[\mathrm{SQ}<0.75],
\label{eq:reward}
\end{equation}
which is the per-step form of Sentinel's own checkpoint-selection score (its
Eq.~(15), $w_{\text{sev}}=0.5$). An operator can raise $w_{\text{sev}}$ to
express a safety-first intent. The advantage is group-relative, with no value
network~\cite{shao2024deepseekmath}:
\begin{equation}
A_{b,g}=\frac{R_{b,g}-\operatorname{mean}_{g'}R_{b,g'}}{\operatorname{std}_{g'}R_{b,g'}+\epsilon},
\label{eq:adv}
\end{equation}
and it is assigned to all $H$ actions of member $g$. Because the members differ
only in their own actions, $A_{b,g}$ is a paired, low-variance comparison. The
traffic noise that dominates a single rollout's return cancels within the
group. We test this directly by ablating CRN, which gives every member an
independent tape.
The policy (the same two-layer, 64-unit Beta/Categorical network as Sentinel,
without a critic) maximises the clipped surrogate
\begin{equation}
\mathcal{J}(\theta)=\mathbb{E}\Big[\min\big(\rho A,\operatorname{clip}(\rho,1\!\pm\!\epsilon)A\big)\Big]
-\beta_{\text{KL}}\,\mathbb{E}\big[D_{\mathrm{KL}}(\pi_\theta\,\|\,\pi_{\text{ref}})\big]
+c_{\text{H}}\,\mathbb{E}[\mathcal{H}(\pi_\theta)],
\label{eq:grpo}
\end{equation}
where $\rho=\pi_\theta(a|s)/\pi_{\theta_{\text{old}}}(a|s)$, the KL is computed in
closed form for the Beta and Categorical heads, and $\pi_{\text{ref}}$ is the
current archive champion (Section~\ref{sec:anchor}), not the initial policy.

\subsection{Solution Archive}
\label{sec:archive}
For each niche $k$ the archive stores the best solution found so far,
$\mathcal{A}[k]$. A solution is any deployable controller together with the
shield parameters it runs under: a neural snapshot, a crystallised tree
program (Section~\ref{sec:crystal}) or a seed program. Seed programs are the
18 constant rules $(\tau,f,d)$ on a coarse grid, which give the archive
interpretable baselines from iteration 0. Every $K$ iterations the current
policy and its crystallised program are \emph{offered} to the archive. For each
niche, the candidate and the incumbent are scored on the same curriculum
scenarios with the same CRN seed, and the incumbent is replaced only if the
candidate is strictly better. Each niche's score on its evaluation set
therefore never decreases. Neural and symbolic solutions compete on equal
terms, so the archive shows directly where a small tree is as good as the
network. It also supplies the teachers for elite injection.

\subsection{Regret-driven scenario curriculum}
\label{sec:curriculum}
Sentinel's attacker maximises the defender's loss. That objective has no
fixed point once attacks exceed what any defender can handle, and its
training shows the late-stage degradation that Sentinel then has to repair
with checkpoint selection. NS-CSD instead ranks scenarios by \emph{regret
against the archive},
\begin{equation}
\mathrm{reg}(\xi)=\max\big(0,\;J(\mathcal{A}[\kappa(\xi)],\xi)-J(\pi_\theta,\xi)\big),
\label{eq:regret}
\end{equation}
which measures how far the live policy falls short of a solution the system
has \emph{already found} for $\xi$. Regret is zero on scenarios that nothing in
the archive handles better, so the curriculum cannot run away into unwinnable
attacks. In UED terms it is a critic-free, archive-grounded replacement for
the value-loss score of prioritised level replay
(PLR)~\cite{jiang2021plr,jiang2021robustplr}. The curriculum keeps a buffer of
64 scenarios. Contexts are replayed from it with probability $0.5$, using
rank-based priority on regret mixed with staleness; the rest are drawn from
the prior. At each discovery step, 16 candidates are proposed: half are
ACCEL-style edits of buffered scenarios~\cite{parkerholder2022accel} and half
are fresh draws. The candidates are scored, and the buffer keeps the
highest-regret scenarios.

\subsection{Archive-guided improvement: elite injection, anchor, roll-back}
\label{sec:anchor}
\emph{Elite injection.} When the archive holds an elite for a context's niche,
one of the $G$ group members is driven by that elite, following mixed-policy
GRPO~\cite{yan2025luffy}. It competes inside the group-relative normalisation
of Eq.~\eqref{eq:adv}. Its actions enter the loss only through an
advantage-weighted likelihood term $-c_{\text{el}}\,A^{+}\log\pi_\theta(a^{\text{el}}|s)$,
so the policy imitates the elite \emph{only where the elite beat the policy's
own samples on identical traffic}.
\emph{Anchor and roll-back.} A validation battery (benign, mild, strong,
chaos; ICMP is excluded so that it stays unseen) is run every $K$ iterations.
Its composite score determines the \emph{champion}, and $\pi_{\text{ref}}$ in
Eq.~\eqref{eq:grpo} is set to it. If validation stays more than a margin below
the champion for three consecutive evaluations, the policy and its shield are
rolled back to the champion. Sentinel selects a checkpoint after training
ends. NS-CSD instead makes the best-known solution the anchor for further
optimisation during training.

\subsection{Certificate-gated shield discovery}
\label{sec:shielddisc}
A diagnostic in Section~\ref{sec:diagnostics} shows that the \emph{fixed}
shield, not the neural policy, limits mitigation of low-rate floods. When link
utilisation is below 0.6, PanicGuard caps the drop rate at 0.1 and raises the
threshold to at least 0.4, so a flood that keeps utilisation under 0.6 passes
almost untouched. NS-CSD therefore treats the PanicGuard parameters
$\phi=(u_{\text{gate}},d_{\text{cap}},\tau_{\text{floor}})$ as part of the
solution and discovers them continuously, under two safeguards.

\emph{Shield probes.} A policy trained behind a fixed clamp gets no learning
signal inside the clamped region: every sampled action is repaired to the same
action and, under CRN, earns the same return. NS-CSD therefore runs two members
of every group under a perturbed shield $\phi+\varepsilon$. This exposes the
policy to the relaxed region and yields a group-relative evolution-strategy
direction $\sum A_{b,g}\varepsilon_{b,g}$ for $\phi$. For the same reason, NS-CSD
does not use Sentinel's shield-repair penalty, which trains the policy to
propose only actions the current shield already allows.

\emph{Safety certificate.} At each discovery step, candidate shields are
proposed along the ES direction (three step sizes) and at random. A candidate
is adopted only if it (a) improves the validation composite by more than a
margin and (b) causes \emph{no safety regression relative to the
operator-approved default shield} on a certification battery of attack-free
traffic, with and without flash crowds. Regression means lower service quality
(tolerance $2\times10^{-3}$ and $5\times10^{-3}$) or any additional severe or
degraded step. The certificate is empirical, based on fixed common random
numbers; it is not a formal proof. The ICMP rules and the remaining guards are
never modified.

\subsection{Continuous rule crystallisation}
\label{sec:crystal}
At every discovery step the shielded policy is distilled into three depth-4
trees (focus classifier; threshold and drop regressors) fitted on
(observation, deployed action) pairs from the validation battery and from
random prior scenarios, with fidelity measured on a 30\,\% held-out split. The
tree program runs behind the same shield and is offered to the archive as a
solution in its own right. For operators we also report the tree after merging
redundant splits, i.e.\ splits whose sub-trees give the same decision
(Section~\ref{sec:rules}).

\begin{algorithm}[t]
\caption{NS-CSD training loop (one seed)}
\label{alg:nscsd}
\small
\begin{algorithmic}[1]
\STATE init $\pi_\theta$; $\pi_{\text{ref}}\!\leftarrow\!\pi_\theta$; shield $\phi\!\leftarrow\!\phi_0$; archive $\mathcal{A}\!\leftarrow$ seed programs; curriculum buffer $\mathcal{B}\!\sim$ prior
\FOR{$i=1,\dots,I$}
  \STATE draw $B$ scenarios from $\mathcal{B}$ (prob.\ 0.5) or the prior; compile schedules + tapes; warm up $W$ steps
  \STATE fork each context into $G$ CRN members: policy samples, 2 shield probes $\phi+\varepsilon$, 1 elite $\mathcal{A}[\kappa(\xi)]$
  \STATE roll out $H$ steps through $\mathcal{S}_\phi$; returns $R_{b,g}$ (Eq.~\ref{eq:reward}); advantages (Eq.~\ref{eq:adv})
  \STATE update $\theta$ with Eq.~\eqref{eq:grpo} + elite likelihood term; accumulate ES direction for $\phi$
  \IF{$i \bmod K = 0$}
    \STATE validate $\pi_\theta$; propose shields; adopt best certified improving $\phi$ (Sec.~\ref{sec:shielddisc})
    \STATE update champion; $\pi_{\text{ref}}\!\leftarrow$ champion; roll back on sustained regression
    \STATE crystallise $\pi_\theta$ into tree program $T$; offer $\pi_\theta$ and $T$ to $\mathcal{A}$ (monotone per niche)
    \STATE propose 16 scenarios; score regret (Eq.~\ref{eq:regret}); keep top-regret in $\mathcal{B}$
  \ENDIF
\ENDFOR
\STATE \textbf{return} champion (policy + shield), best tree program, archive
\end{algorithmic}
\end{algorithm}
